\documentclass[a4paper]{article}
\usepackage[normalem]{ulem}
\usepackage[font=small,labelfont=bf]{caption}

% impostazioni generali
%Tutti gli usepackage vanno qui
\usepackage{geometry}
\usepackage[italian]{babel}
\usepackage[utf8]{inputenc}
\usepackage{tabularx}
\usepackage{longtable}
\usepackage{hyperref}
\usepackage{enumitem}
\usepackage{array} 
\usepackage{booktabs}
\newcolumntype{M}[1]{>{\centering\arraybackslash}m{#1}}
\usepackage[toc]{appendix}

\hypersetup{
	colorlinks=true,
	linkcolor=blue,
	filecolor=magenta,
	urlcolor=blue,
}
% Numerazione figure
\let\counterwithout\relax
\let\counterwithin\relax
\usepackage{chngcntr}

% distanziare elenco delle figure e delle tabelle
\usepackage{tocbasic}
\DeclareTOCStyleEntry[numwidth=3.5em]{tocline}{figure}% for figure entries
\DeclareTOCStyleEntry[numwidth=3.5em]{tocline}{table}% for table entries


\counterwithin{table}{subsection}
\counterwithin{figure}{subsection}

\usepackage[bottom]{footmisc}
\usepackage{fancyhdr}
\setcounter{secnumdepth}{4}
\usepackage{amsmath, amssymb}
\usepackage{array}
\usepackage{graphicx}

\usepackage{ifthen}

\usepackage{float}
\restylefloat{table}

\usepackage{layouts}
\usepackage{url}
\usepackage{comment}
\usepackage{eurosym}

\usepackage{lastpage}
\usepackage{layouts}
\usepackage{eurosym}

\geometry{a4paper,top=3cm,bottom=4cm,left=2.5cm,right=2.5cm}

%Comandi di impaginazione uguale per tutti i documenti
\pagestyle{fancy}
\lhead{\includegraphics[scale=0.1]{../../../template/images/logo_no_motto.jpeg}}
%Titolo del documento
\rhead{\doctitle{}}
%\rfoot{\thepage}
\cfoot{Pagina \thepage\ di \pageref{LastPage}}
\setlength{\headheight}{35pt}
\setcounter{tocdepth}{5}
\setcounter{secnumdepth}{5}
\renewcommand{\footrulewidth}{0.4pt}

% multirow per tabelle
\usepackage{multirow}

% Permette tabelle su più pagine
%\usepackage{longtable}


% colore di sfondo per le celle
\usepackage[table]{xcolor}

%COMANDI TABELLE
\newcommand{\rowcolorhead}{\rowcolor[HTML]{007c95}}
\newcommand{\cellcolorhead}{\cellcolor[HTML]{007c95}}
\newcommand{\hlinetable}{\arrayrulecolor[HTML]{007c95}\hline}

%intestazione
% check for missing commands
\newcommand{\headertitle}[1]{\textbf{\color{white}#1}} %titolo colonna
\definecolor{pari}{HTML}{b1dae3}
\definecolor{dispari}{HTML}{d7f2f7}

% comandi glossario
\newcommand{\glo}{$_{G}$}
\newcommand{\glosp}{$_{G}$ }


%label custom
\makeatletter
\newcommand{\uclabel}[2]{%
	\protected@write \@auxout {}{\string \newlabel {#1}{{#2}{\thepage}{#2}{#1}{}} }%
	\hypertarget{#1}{#2}
}
\makeatother

%riportare pezzi di codice
\definecolor{codegray}{gray}{0.9}
\newcommand{\code}[1]{\colorbox{codegray}{\texttt{#1}}}



% dati relativi alla prima pagina
% Configurazione della pagina iniziale
\newcommand{\doctitle}{\textit{Verbale} interno 01-12-2022 }
\newcommand{\docdate}{01 Dicembre 2022}
\newcommand{\rev}{1.0.0}
\newcommand{\stato}{Approvato}
\newcommand{\uso}{Interno}
\newcommand{\approv}{Elena Pandolfo}
\newcommand{\red}{Tommaso Allegretti}
\newcommand{\ver}{Gabriele Mantoan\\& Mirko Stella}
\newcommand{\dest}{\textit{Seven Clickers}
									  \\ Prof. Vardanega Tullio 
									   \\ Prof. Cardin Riccardo}
\newcommand{\describedoc}{\textit{Verbale} riguardante il meeting tenuto il giorno 01-12-2022}


 % editare questo

\makeindex

%comando per far andare a capo i paragrafi
\makeatletter
\renewcommand\paragraph{
\@startsection {paragraph}{4}{0mm}{-\baselineskip}{.5\baselineskip}{\normalfont \normalsize \bfseries }}
\makeatother

\newcommand*{\fullref}[1]{\hyperref[{#1}]{\ref*{#1} \nameref*{#1}}}

\begin{document}
\counterwithin{table}{section}

% Prima pagina
\thispagestyle{empty}
\renewcommand{\arraystretch}{1.3}

\begin{titlepage}
	\begin{center}
		
	\includegraphics[scale = 0.40]{../../../template/images/logo.jpeg}
	\\[1cm]
	\href{mailto:7clickersgroup@gmail.com}		      	
	{\large{\textit{7clickersgroup@gmail.com} } }\\[2.5cm]
	\Huge \textbf{\doctitle} \\[1cm]
	 \large
			 \begin{tabular}{r|l}
                        \textbf{Versione} & \rev{} \\
                        \textbf{Stato} & \stato{} \\
                        \textbf{Uso} & \uso{} \\                         
                        \textbf{Approvazione} & \approv{} \\                      
                        \textbf{Redazione} & \red{} \\ 
                        \textbf{Verifica} &  \ver{} \\                         
                        \textbf{Distribuzione} & \parbox[t]{5cm}{ \dest{} }
                \end{tabular} 
                \\[3.5cm]
                \large \textbf{Descrizione} \\ \describedoc{} 
     \end{center}
\end{titlepage}

% Diario delle modifiche
\section*{Registro delle modifiche}

\newcommand{\changelogTable}[1]{
	 

\renewcommand{\arraystretch}{1.5}
\rowcolors{2}{pari}{dispari}
\begin{longtable}{ 
		>{\centering}M{0.07\textwidth} 
		>{\centering}M{0.13\textwidth}
		>{\centering}M{0.20\textwidth}
		>{\centering}M{0.17\textwidth} 
		>{\centering\arraybackslash}M{0.30\textwidth} 
		 }
	\rowcolorhead
	\headertitle{Vers.} &
	\centering \headertitle{Data} &	
	\headertitle{Autore} &
	\headertitle{Ruolo} & 
	\headertitle{Descrizione} 
	\endfirsthead	
	\endhead
	
	#1

\end{longtable}
\vspace{-2em}

}



\changelogTable{
	1.0.0 & 14-12-22 & Elena Pandolfo & Responsabile &  Approvazione documento\tabularnewline
	0.1.0 & 01-12-22 & Gabriele Mantoan \\ Mirko Stella & Verificatori & Verifica documento\tabularnewline
	0.0.1 & 09-11-22 & Rino Sincic & Analista &  Redazione documento\\
} % editare questo
\pagebreak

% Indice
{
    \hypersetup{linkcolor=black}
    \tableofcontents
    \listoffigures %elenco figure
}
\pagebreak

% Contenuto
\section{Introduzione}

\subsection{Scopo del documento}
\textit{Norme di Progetto} è il documento che definisce le regole e gli standard che devono essere seguiti durante il ciclo di vita\textsubscript{g} del prodotto\textsubscript{g}.
\subsection{Scopo del prodotto\textsubscript{g}}
Il prodotto\textsubscript{g} in questione nasce dalla necessità dell'azienda SanMarco Informatica di fornire una soluzione agli sprechi
derivati dall'adozione di uno ShowRoom tradizionale proponendo uno ShowRoom3D che sia ugualmente o ancora più immersivo.
\subsection{Glossario}
In questo documento sono state segnate con il pedice "g" tutte le parole che, secondo noi, necessitano di una loro definizione più accurata nel documento di \textit{Glossario v1.0.0}.
\subsection{Riferimenti}
\subsubsection{Riferimenti normativi}
\begin{itemize}
    \item Regolamento del Progetto Didattico: \url{https://www.math.unipd.it/~tullio/IS-1/2022/Dispense/PD02.pdf}
    \item Capitolato di appalto C6 - ShowRoom3D: \url{https://www.math.unipd.it/~tullio/IS-1/2022/Progetto/C6.pdf}
\end{itemize}

\subsubsection{Riferimenti informativi}
\begin{itemize}
	\item Materiale didattico Ingegneria del Software - T02 Processi di ciclo di vita\textsubscript{g}: \url{https://www.math.unipd.it/~tullio/IS-1/2022/Dispense/T02.pdf}
	\item Materiale didattico Ingegneria del Software - T03 Il ciclo di vita del SW: \url{https://www.math.unipd.it/~tullio/IS-1/2022/Dispense/T03.pdf}
	\item Materiale didattico Ingegneria del Software - T04 Gestione di progetto: \url{https://www.math.unipd.it/~tullio/IS-1/2022/Dispense/T04.pdf}
	\item Materiale didattico Ingegneria del Software - T06 Analisi dei requisiti: \url{https://www.math.unipd.it/~tullio/IS-1/2022/Dispense/T06.pdf}
	\item Materiale didattico Ingegneria del Software - T010 Verifica e validazione: introduzione \url{https://www.math.unipd.it/~tullio/IS-1/2022/Dispense/T010.pdf}
	\item Lo standard ISO/IEC 12207:1995 : \url{https://www.math.unipd.it/~tullio/IS-1/2009/Approfondimenti/ISO_12207-1995.pdf}
\end{itemize}
\pagebreak

\section{Processi Primari}
I processi primari devono essere seguiti durante il ciclo di vita\textsubscript{g} del software.
Nel nostro il ciclo di vita\textsubscript{g} del software non comprende l'installazione e la manutenzione ma termina con lo sviluppo.
Il progetto in questione è da considerarsi un progetto didattico.
\subsection{Fornitura}

\subsubsection{Scopo}
Lo scopo del processo di fornitura è quello di stabilire quali sono le attività, 
le risorse e i compiti necessari allo svolgimento del progetto. 

\subsubsection{Descrizione}
Nella seguente sezione sono riportate le norme che il gruppo si impegna a rispettare al fine di creare e 
mantenere un rapporto attivo e vantaggioso con il proponente SanMarco Informatica. 

\subsubsection{Rapporti con il proponente}
Il gruppo si manterrà in contatto con il proponente regolarmente, organizzando incontri su Google Meet o 
comunicando in modo asincrono con email o Google Chat. Di seguito un elenco dei principali temi che possono essere 
motivo di confronto con il proponente: 
\begin{itemize} 
	\item Chiarimento di eventuali dubbi o incomprensioni; 
	\item Riscontro sul lavoro svolto;
	\item Definizione dei requisiti obbligatori e opzionali;
	\item Opinioni riguardanti tecnologie da utilizzare;
	\item Proposte da parte del gruppo di soluzioni alternative.
\end{itemize}
\subsection{Sviluppo}

\subsubsection{Scopo}
Lo scopo del processo di sviluppo è quello di stabilire le attività necessarie allo sviluppo del prodotto software.

\subsubsection{Descrizione}
Nelle seguenti sezioni verranno descritte più nel dettaglio le principali attività, relative al processo di sviluppo:
\begin{itemize} 
    \item Analisi dei requisiti;
	\item Progettazione;
	\item Codifica.
\end{itemize}

\subsubsection{\textit{Analisi dei Requisiti}}
L’\textit{Analisi dei Requisiti} è un’attività fondamentale che precede la progettazione. 
Si concretizza nel documento \textit{Analisi dei Requisiti}. Lo scopo di quest’attività è:
    \begin{itemize}
        \item Definire le funzionalità che il prodotto\textsubscript{g} andrà ad offrire;
        \item Porre le basi per la fase di progettazione del software;
        \item Fare una stima della mole di lavoro;
        \item Facilitare la verifica\textsubscript{g}. 
    \end{itemize}

\paragraph{Struttura del documento} Il documento è strutturato nel modo seguente:
    \begin{itemize}
        \item \textbf{Introduzione}: contiene una breve descrizione del documento e specifica gli attori dei casi d’uso;
		\item \textbf{Casi d’uso}: vengono specificati i casi d’uso individuati;
		\item \textbf{Requisiti}: sono classificati i requisiti che il prodotto\textsubscript{g} finale dovrà soddisfare.
    \end{itemize}

\paragraph{Casi d’uso} 

I casi d’uso sono identificati seguendo la seguente struttura:

\begin{center}\textbf{UC [Numero caso d’uso].[Numero sottocaso d’uso]-[Titolo caso d’uso]}\end{center}

\noindent Vengono poi indicati:
    \begin{itemize}
        \item \textbf{Diagramma UML\textsubscript{g}}: viene mostrato il diagramma solo per i casi più complessi, in cui è ritenuto necessario per una maggiore comprensione. Per i casi d’uso facoltativi il diagramma è colorato di azzurro (non standard UML\textsubscript{g}) per una maggiore differenziazione;
	    \item \textbf{Attore\textsubscript{g} primario}: utente esterno al sistema che svolge il caso d’uso;
	    \item \textbf{Descrizione}: breve descrizione del caso d’uso;
	    \item \textbf{Precondizioni}: indica la condizione del sistema prima del verificarsi del caso d’uso;
	    \item \textbf{Postcondizioni}: indica la condizione del sistema dopo che si è verificato il caso d’uso;
	    \item \textbf{Scenario principale}: descrive l’interazione tra l’attore\textsubscript{g} primario e il sistema;
	    \item \textbf{Estensioni}: casi d’uso alternativi che possono verificarsi al posto di quello a cui sono collegati. 
    \end{itemize}

\paragraph{Requisiti} 
I requisiti sono classificati secondo la seguente struttura: 

\begin{center}\textbf{R.[Tipologia][Numero seriale]}\end{center}

\noindent con:
    \begin{itemize}
        \item \textbf{Tipologia}: 
            \begin{itemize}
                \item \textbf{F}: funzionale;
                \item \textbf{Q}: qualitativo;
                \item \textbf{D}: di dominio;
                \item \textbf{P}: prestazionale.
            \end{itemize}
	    \item \textbf{Descrizione}: breve descrizione del requisito;
	    \item \textbf{Classificazione}: un requisito può essere facoltativo o obbligatorio;
	    \item \textbf{Fonti}: possono essere: 
            \begin{itemize}
                \item Casi d’uso;
                \item Capitolato;
                \item Decisione interna.
            \end{itemize}
    \end{itemize}

\subsubsection{Progettazione} 
L’attività di progettazione è assegnata ai progettisti. Seguendo quanto indicato 
nell’\textit{Analisi dei Requisiti}, l’obiettivo è quello di progettare l’architettura del sistema, 
prima con la creazione di un Proof of Concept\textsubscript{g} per la Requirments and Tecnology Baseline e poi andando 
nel dettaglio per la Product Baseline.

\paragraph{Requirments and Tecnology Baseline} 
In questa fase vengono fissati i requisiti che il gruppo si impegna a soddisfare, in accordo con il 
proponente\textsubscript{g}; si studiano le tecnologie, i framework\textsubscript{g} e le librerie utili alla realizzazione del 
prodotto\textsubscript{g} finale e si crea di conseguenza un Proof of Concept\textsubscript{g}. I materiali da mostrare sono:
    \begin{itemize}
        \item Tecnologie, framework\textsubscript{g}, librerie utilizzate, motivando la scelta;
        \item Proof of Concept\textsubscript{g};
        \item \textit{Analisi dei Requisiti}.
    \end{itemize}
La documentazione da mostrare:
    \begin{itemize}
        \item \textit{Piano di Progetto};
        \item \textit{Piano di Qualifica};
        \item \textit{Norme di Progetto};
        \item \textit{Verbali} interni ed esterni.
    \end{itemize}


\pagebreak
\section{Processi di Supporto}
I processi di supporto supportano quelli primari in modo da renderli più efficienti ed efficaci.
\section{Documentazione}
Questa sezione descrive le convenzioni,gli strumenti e le modalità con cui il gruppo si impegna a stilare la documentazione interna ed esterna relativa al progetto.

\subsubsection{Convenzioni generali}
Le convenzioni di seguito riportate vengono applicate a tutti i documenti.
Esse rendono i documenti stilati omogenei tra loro contribuendo a rendere il progetto professionale.

\paragraph{Versionamento}
Il numero di versione permette di capire lo stato in cui si trova un documento.
Un documento può trovarsi nei seguenti stati:
\begin{itemize} 
    \item \textbf{Approvato}: il documento è verificato ed approvato dal Responsabile;
    \item \textbf{Verificato}: il documento risulta verificato ma non ancora visionato dal Responsabile;
    \item \textbf{In Sviluppo}: sono presenti delle modifiche che non sono state verificate.
\end{itemize}
Il numero di versione ha il formato \textbf{X.Y.Z} dove:
\begin{itemize} 
    \item \textbf{X}: indica una versione approvata dal Responsabile, la numerazione parte da 0
    e la prima versione approvata è la 1.0.0;
    \item \textbf{Y}: indica una versione verificata dal Verificatore, la numerazione inizia da 0 e si azzera ad ogni incremento di X. La prima versione
    verificata è la 0.1.0;
    \item \textbf{Z}: indica una versione in fase di modifica da parte dei redattori che ne incrementano il numero ad ogni modifica,
    la numerazione parte da 1 e si azzera ad ogni incremento di X o Y. La prima versione modificata è la 0.0.1.
\end{itemize}

\paragraph{Struttura Generale}
Ogni documento deve presentare le seguenti sezioni nell'ordine in cui vengono presentate:
\begin{itemize} 
    \item \textbf{Intestazione}:
    contiene:
    \begin{itemize} 
        \item Logo compreso di motto;
        \item Indirizzo email di gruppo; 
        \item Titolo;
        \item Tabella contenente le informazioni generali:
        \begin{itemize}
            \item Versione;
            \item Stato;
            \item Uso;
            \item Approvazione: indica il responsabile di progetto che ha approvato il documento; 
            \item Redazione: elenco dei collaboratori che hanno partecipato alla stesura del documento;
            \item Verifica: elenco dei verificatori che hanno verificato il documento;
            \item Distribuzione: elenco delle persone o organizzazioni a cui è destinato il documento.
        \end{itemize}
        \item Breve descrizione del documento .
    \end{itemize}
    \item \textbf{Registro delle modifiche}:
    tabella che identifica ogni versione del documento indicandone:
    \begin{itemize} 
        \item Versione;
        \item Data;
        \item Autore;
        \item Ruolo;
        \item Descrizione.
    \end{itemize}
    \item \textbf{Indice}:
    elenco ordinato dei titoli dei capitoli, ovvero delle varie parti di cui si compone il documento;
    \item \textbf{Contenuto}:
    varia a seconda del tipo di documento.
\end{itemize}
\paragraph{Date}
Le date devono rispettare il seguente formato: \textbf{dd-mm-yyyy}
All'interno delle tabelle il formato deve essere il seguente: \textbf{dd-mm-yy}
\paragraph{Nomi di persona}
All'interno dei documenti i nomi di persona rispetteranno l'ordine nome seguito dal cognome della persona menzionata.

\paragraph{Elenchi puntati}
Gli elenchi puntati devono rispettare le seguenti regole:
\begin{itemize} 
    \item Ogni elemento dell'elenco deve iniziare con la lettera maiuscola;
    \item Ogni elemento dell'elenco deve terminare con ";" ad eccezione dell'ultimo elemento
    che deve terminare con "."; 
    \item Dopo i due punti la frase deve iniziare con la lettera minuscola.
\end{itemize}

\paragraph{Stile del testo}
\begin{itemize} 
    \item \textbf{Grassetto}: stile utilizzato per i titoli delle sezioni e per i primi termini degli elenchi puntati;
    \item \textbf{Corsivo}: viene utilizzato per citare il nome dei documenti, ad esempio \textit{Piano di Progetto}. 
\end{itemize}

\subsection{Verifica}
La verifica viene svolta da due verificatori prima del merge con il branch documentation.
Consiste nell'esaminare i file prodotti da chi ne ha fatto la stesura e segnalarne la non validità o 
la presenza di errori nei concetti esposti.\\
Un verificatore dovrà verificare il documento a partire dalle modifiche fatte dopo l'ultima versione verificata.
Le modifiche da verificare quindi possono essere dedotte dal registro dei cambiamenti presente in ogni documento.
Una volta controllato il documento, il primo verificatore segnalerà eventuali errori e
successivamente dovrà spuntare come approvata la Pull Request nella sezione dedicata su GitHub\textsubscript{g}. \\
A questo punto se un secondo verificatore noterà la necessità di qualche altro cambiamento da apportare, chi dovrà apportare le modifiche farà una pull in locale per
allineare il proprio branch con quello in remoto e continuare con il proprio lavoro.\\
Dopo il push delle modifiche se i file risultano corretti anche dal secondo verificatore esso aggiungerà i nomi dei verificatori all'intestazione modificando il file titlepage\_input.tex
 e creerà la riga nel registro delle modifiche, nel file changelog\_input.tex,  inserendo la nuova versione secondo le regole di versionamento e scrivendo nella colonna Autore sia il suo nome,
che quello del primo verificatore, e in descrizione "Verifica".
Dopo aver eseguito un commit sul ramo da integrare approva la Pull Request e conferma il merge secondo le norme di progetto descritte nella sezione dedicata.
\subsection{Approvazione}
L’approvazione viene svolta dal Responsabile di Progetto.\\
L'approvazione consiste nell'aprire una Pull Request di approvazione. Nel caso in cui il Responsabile di progetto riscontrasse ulteriori problematiche, segnalerà ai Verificatori le eventuali modifiche da apportare. I verificatori apporteranno tali modifiche prima di chiudere la Pull Request di approvazione e quindi di integrare i cambiamenti nel branch documentation.\\
Per i verbali si effettua questa pratica non appena il file prodotto viene verificato e quindi passa tutte le eventuali reviews create dai verificatori mentre per tutti gli altri documenti prima di una consegna.\\
Ad esito positivo di approvazione, il Responsabile di Progetto creerà una nuova riga e compilerà i campi nelle colonne corrispondenti del registro delle modifiche inserendo: l'ultima versione secondo le norme di versionamento, la data, il proprio nome, il suo ruolo e la voce "Approvazione" nell'ultima colonna.
\subsubsection{Strumenti per la stesura}
\begin{itemize} 
    \item LaTeX: è un linguaggio di marcatura per la preparazione di testi, basato sul 				  		  programma di composizione tipografica TEX.\\
	Nel branch documentation  si possono trovare i file .pdf prodotti e la cartella “latex”. La cartella latex contiene tre cartelle interne:
	\begin{itemize}
	\item esterni e interni, contengono file .tex di documentazione esterna ed interna come ad esempio i verbali o altra documentazione esterna/interna:
	\begin{itemize}
		\item la cartella config, contiene i file .tex con le parti fisse dei documenti (intestazione,registro delle modifiche,tracciamento dei temi affrontati) che vengono modificati con i dati del documento specifico 
		\item la cartella res/sections, contiene i file .tex con il contenuto vero e proprio (sezioni del documento) che viene redatto in maniera libera dal redattore
		\item un file col nome del documento pdf con estensione .tex che viene compilato per produrre il file pdf
	\end{itemize}
	 \item template, contiene file .tex di base utilizzati secondo necessità per comporre i documenti:
	 \begin{itemize}
	 \item changelox.tex è il file di template che serve per scrivere il registro delle modifiche
	 \item package.tex è il file che contiene tutti gli usepackage\textsubscript{g}
	 \item titlepage.tex è il file di template che contiene la configurazione della pagina iniziale di ogni documento  
	 \item tracking.tex è il file di template che contiene il tracciamento dei temi affrontati nel documento
	 \end{itemize}
	\end{itemize}
\end{itemize}

\subsection{Gestione della configurazione}
\subsubsection{Scopo}
La gestione della configurazione è un processo\textsubscript{g} che mira a gestire e controllare i cambiamenti apportati a un prodotto\textsubscript{g} software o a un sistema durante il suo ciclo di vita\textsubscript{g}. 
La gestione della configurazione per la documentazione descrive come vengono identificate, controllate, tracciate e gestite le versioni di un documento.
\subsubsection{Descrizione}
In questa sezione sono riportate tutte le norme relative alla configurazione degli strumenti utilizzati per il tracciamento e l’organizzazione dei documenti e del codice.
\subsubsection{Versionamento}
Il numero di versione permette di capire lo stato in cui si trova un documento.
Un documento può trovarsi nei seguenti stati:
\begin{itemize} 
    \item \textbf{Approvato}: il documento è verificato ed approvato dal responsabile;
    \item \textbf{Verificato}: il documento risulta verificato ma non ancora visionato dal responsabile;
    \item \textbf{In Sviluppo}: sono presenti delle modifiche che non sono state verificate.
\end{itemize}
Il numero di versione ha il formato \textbf{X.Y.Z} dove:
\begin{itemize} 
    \item \textbf{X}: indica una versione approvata dal responsabile, la numerazione parte da 0
    e la prima versione approvata è la 1.0.0;
    \item \textbf{Y}: indica una versione verificata dal verificatore, la numerazione inizia da 0 e si azzera ad ogni incremento di X. La prima versione
    verificata è la 0.1.0;
    \item \textbf{Z}: indica una versione in fase di modifica da parte dei redattori che ne incrementano il numero ad ogni modifica,
    la numerazione parte da 1 e si azzera ad ogni incremento di X o Y. La prima versione modificata è la 0.0.1.
\end{itemize}

\paragraph{Sistemi software utilizzati}
Per il versionamento il gruppo ha deciso di utilizzare un repository\textsubscript{g} GitHub\textsubscript{g}, un servizio di hosting per progetti software che implementa uno strumento di controllo versione distribuito Git\textsubscript{g}.
Per informazioni più approfondite sulla struttura e la gestione del repository\textsubscript{g} si veda sezione \fullref{GitHub}.
\subsection{Accertamento della qualità}

\subsubsection{Scopo}
Lo scopo del processo\textsubscript{g} di accertamento della qualità è quello di garantire e rispettare i requisiti di qualità preposti.
\subsubsection{Descrizione}
Questo processo\textsubscript{g} è direttamente collegato al documento di \textit{Piano di Qualifica} in cui il gruppo si impegna a mantenere, con l'adozione di metriche concordate e discusse, la qualità nei processi e nei prodotti.
\subsubsection{Attività per il controllo di qualità}
Queste sono le attività che ogni membro deve rispettare per il mantenimento della qualità.
Ogni membro del gruppo deve:
\begin{itemize}
\item Comprendere le attività da svolgere e gli obiettivi da raggiungere;
\item Aver compreso le metriche inserite nel documento di \textit{Piano di Qualifica} e rispettare le normative e gli standard definiti al suo interno;
\item Mantenere le normative inserite nelle \textit{Norme di Progetto};
\item Riconoscere nel documento su cui si sta lavorando valori di cui tenere conto per un'analisi successiva;
\item Utilizzare il sistema di tracciamento delle issue\textsubscript{g} fornito da Github\textsubscript{g} come descritto nel documento di \textit{Norme di Progetto};
\item Attuare miglioramento continuo, ponendosi obiettivi incrementali.
\end{itemize}
\paragraph{Attuazione del controllo di qualità}
Il documento deve essere aggiornato ad ogni aggiunta nel numero di versione centrale ,come descritto nella sezione di \ref{versionamento} Versionamento , in cui si andrà ad aggiornare i grafici e le misure scritte nel documento. L’analisi e il commento di ogni misura inserita nei grafici verrà successivamente fatta in fase di retrospezione dello Sprint.
\subsubsection{Denominazione requisiti nel \textit{Piano di Qualifica}}
 Per la denominazione nel tracciamento dei requisiti nella sezione di \textit{Specifica dei Test} nel documento di \textit{Piano di Qualifica} viene adottata la seguente scrittura:
 \begin{center}\textbf{T[Tipologia di Test]R[Tipologia di Requisito][Numero Identificativo]}\end{center}
 [Tipologia di Test] indica a che tipo di Test si sta riferendo:
 \begin{itemize}
 \item Test di Unità corrisponde a "TU";
 \item Test di Integrità corrisponde a "TI";
 \item Test di Accettazione "TA";
 \item Test di Sistema corrisponde a "TS".
 \end{itemize}
 [Tipologia di Requisito] indica a che tipo di Requisito si sta riferendo:
 \begin{itemize}
 \item Requisito Funzionale corrisponderà a "RF";
 \item Requisito Qualitativo corrisponderà a "RQ";
 \item Requisito di Dominio corrisponderà a "RD";
 \item Requisito Prestazionale corrisponderà a "RP".
 \end{itemize}
 [Numero Indicativo] indica il numero del Requisito descritto nel documento di \textit{Analisi dei Requisiti}.
 
\subsubsection{Denominazione degli obiettivi di qualità nel \textit{Piano di Qualifica}}
Per la denominazione delle metriche di prodotto\textsubscript{g} è stato adottata la seguente scrittura:
\begin{center}\textbf{MPD[Numero]}\end{center}
Il Numero indicato è un numero che incrementa in base al numero di metriche per sezione.
Per la denominazione delle metriche di processo\textsubscript{g} è stato adottata la seguente scrittura:
\begin{center}\textbf{MPC[Numero]}\end{center}
Il Numero indicato è un numero che incremente in base al numero di metriche per sezione.
Viene inoltre inserito il nome della metrica, il valore minimo e il valore ottimo.\\
La descrizione dettagliata di ogni metrica la si può trovare alla sezione \fullref{Metriche_qualita}.
\subsubsection{Metriche} 
Per tracciare la soddisfazione dei requisiti è stata concordata una semplice metrica.
\begin{itemize}
    \item \textbf{MPC11: Metriche soddisfatte}.
\end{itemize}

\input{res/sections/processi_supporto/verifica\textsubscript{g}/verifica_main.tex}
\input{res/sections/processi_supporto/validazione\textsubscript{g}.tex}

\pagebreak
\section{Processi Organizzativi}
I processi organizzativi mirano a gestire i processi e il loro miglioramento, l'organizzazione degli strumenti 
di supporto e la gestione del personale.
\subsection{Gestione Organizzativa}
\subsubsection{Scopo}
Questo processo ha come scopo quello di gestire al meglio i vari processi. Contiene tutte le attività e i compiti utili a tale scopo.
\subsubsection{Descrizione}
In questa sezione sono descritte tutte le norme che servono per gestire al meglio l’organizzazione del gruppo, le comunicazioni, la ripartizione dei compiti e dei ruoli e le attività da svolgere.

\subsubsection{Ruoli}
I componenti del gruppo si suddivideranno nei seguenti ruoli per periodi di circa 2-3 settimane (dipendentemente dalle esigenze del periodo) e al termine del periodo i ruoli verranno risuddivisi. 
Visto che nelle varie fasi di sviluppo del progetto le attività da svolgere variano, non sempre sarà necessario coprire tutti i ruoli.\\
Inoltre sarà necessario tenere traccia delle ore che ogni componente dedica al progetto ed il ruolo associato a quelle ore, in modo da andare a rispettare la tabella degli impegni individuali.\\
I ruoli e le loro competenze sono i seguenti:

\paragraph{Responsabile}
Deve avere la visione d'insieme del progetto e coordinare i membri, inoltre si occupa di rappresentare il gruppo con le interazione esterne (proponente, committente ecc...). Le sue competenze specifiche sono:
\begin{itemize}
	\item Ad ogni iterazione\textsubscript{g} c'è un solo responsabile;
	\item Presentare il diario di bordo in aula;
	\item Redarre l'ordine del giorno prima di ogni meeting interno del gruppo;
	\item Suddivide le attività del gruppo in singole issue (ma non le assegna ai membri del gruppo);
	\item In fase di release\textsubscript{g} si occupa di approvare\textsubscript{g} tutti i documenti che necessitano approvazione.
\end{itemize}

\paragraph{Analista}
Si occupa di trasformare i bisogni del proponente nelle aspettative che il gruppo deve soddisfare per sviluppare un prodotto professionale. Le sue competenze specifiche sono:
\begin{itemize}
	\item Interrogare il proponente riguardo allo scopo del prodotto e le funzionalità che deve avere;
	\item Studiare le risposte del proponente per identificare i requisiti\textsubscript{g} e redarre l' \textit{Analisi dei Requisiti}.
\end{itemize}

\paragraph{Amministratore}
Si occupa del funzionamento, mantenimento e sviluppo degli strumenti tecnologici usati dal gruppo. Le sue competenze specifiche sono:
\begin{itemize}
	\item Ad ogni iterazione\textsubscript{g} basta un solo amministratore;
	\item Gestione delle segnalazioni e problemi dei membri del gruppo riguardanti problemi e malfunzionamenti con gli strumenti tecnologici;
	\item Valuta l'utilizzo di nuove tecnologie e ne fa uno studio preliminare per poter presentare al gruppo i pro e i contro del suo utilizzo.
\end{itemize}

\paragraph{Progettista}
Si occupa di scegliere la modalità migliore per soddisfare le aspettative del committente che gli analisti hanno ricavato dall'analisi dei requisiti\textsubscript{g}. Le sue competenze specifiche sono:
\begin{itemize}
	\item Scegliere eventuali pattern architetturali da implementare;
	\item Sviluppare lo schema UML\textsubscript{g} delle classi\textsubscript{g}.
\end{itemize}

\paragraph{Programmatore}
Si occupa di implementare le scelte e i modelli fatti dal progettista. Le sue competenze specifiche sono:
\begin{itemize}
	\item Scrivere il codice atto a implementare lo schema delle classi;
	\item Scrivere eventuali test;
	\item Scrivere la documentazione per la comprensione del codice che scrive.
\end{itemize}

\paragraph{Verificatore}
Si occupa di controllare che ogni file che viene caricato in un branch protetto della repository\textsubscript{g} sia conforme alle norme di progetto. Le sue competenze specifiche sono:
\begin{itemize}
	\item Controllare i file modificati o aggiunti durante una pull request tra un ramo non protetto e un ramo protetto siano conformi alle norme di progetto e cercano errori di altra natura (ortografici, sintattici, logici, build ecc...).
\end{itemize}
\subsubsection{Gestione delle comunicazioni}
Le comunicazioni per tutta la durata del progetto potranno essere interne, ovvero tra i componenti del gruppo;
oppure esterne, quindi tra il gruppo e una o più persone esterne al gruppo, in genere il proponente\textsubscript{g} o il committente\textsubscript{g}.

\paragraph{Comunicazioni Interne}
Si svolgono tra i componenti del gruppo tramite diversi canali di comunicazione. 
Ogni canale è dedito a un diverso tipo di comunicazione e comporta delle regole da seguire, per aumentare la 
produttività del gruppo.

\begin{itemize}
    \item \textbf{Telegram}: viene utilizzato per delle comunicazioni brevi;
    \item \textbf{Discord\textsubscript{g}}: sono stati creati diversi canali, che possono essere:
    \begin{itemize}
        \item \textbf{Testuali}: questi si suddividono in canali per argomento o per ruolo. I canali per argomento vengono utilizzati 
        esclusivamente per trattare dell'argomento in questione, mentre i canali per ruolo sono utili per le comunicazioni tra i membri del gruppo che
        sono assegnati allo stesso ruolo nello stesso periodo e hanno bisogno di comunicare tra loro;
        \item \textbf{Vocali}: vengono utilizzati per riunioni brevi, principalmente di aggiornamento.
    \end{itemize}
    \item \textbf{Zoom\textsubscript{g}}: viene utilizzato per le riunioni interne del gruppo, al termine delle quali verrà redatto un \textit{Verbale}.
\end{itemize}

\paragraph{Comunicazioni Esterne}
Le comunicazioni esterne si svolgono tra il gruppo e una persona esterna, generalmente il proponente\textsubscript{g} o il committente\textsubscript{g}. Con il 
committente\textsubscript{g} verranno utilizzati i canali suggeriti da quest'ultimo, mentre con il proponente\textsubscript{g} si è stabilito di utilizzare:
 \begin{itemize}
    \item \textbf{e-mail}: per comunicazioni brevi e immediate, si farà uso dell'indirizzo e-mail del gruppo. Il compito di scrivere le e-mail è assegnato al 
    responsabile, che, dopo una breve verifica\textsubscript{g} del messaggio con il gruppo avrà il compito di inviare la e-mail e notificare il gruppo alla risposta;
    \item \textbf{Google Chat}: per messaggi brevi, sempre scritti dal responsabile;
    \item \textbf{Google Meet}: per riunioni esterne in videochiamata. Queste riunioni potranno essere 
    richieste sia dal gruppo, che dal proponente\textsubscript{g} e al termine verrà redatto un \textit{Verbale}. 
\end{itemize}

\subsubsection{Gestione delle riunioni}
Le riunioni si dividono in interne ed esterne. Per ogni riunione verrà redatto un \textit{Verbale}, per poter tenere traccia 
degli argomenti trattati e delle decisioni prese.

\paragraph{Riunioni Interne}
Una riunione interna si svolge esclusivamente tra i membri del gruppo utilizzando il canale apposito Zoom\textsubscript{g}. Per ogni riunione 
il responsabile sarà incaricato di 
\begin{itemize}
    \item Preparare una scaletta degli argomenti da trattare, che potranno essere poi integrati da 
    eventuali punti di discussione portati dagli altri membri del gruppo;
    \item Fare da moderatore e dirigere la riunione per ottimizzare il tempo.
\end{itemize}

Le riunioni si svolgeranno a cadenza settimanale, cercando di trovare giorni e orari agevoli a tutti i membri del gruppo.

\paragraph{Riunioni Esterne}
Le riunioni esterne si svolgono tra i membri del gruppo e una persona esterna. Possono essere richieste da entrambe le parti 
e tramite i canali di comunicazione stabiliti, (Google Meet con il proponente\textsubscript{g}). Se necessario verrà preparata una presentazione 
da mostrare ai soggetti esterni per ottimizzare i tempi e facilitare la comprensione. Al termine della riunione verrà redatto un 
\textit{Verbale}. 

\subsubsection{Metriche}
Abbiamo considerato questa metrica come rilevante per mantenere la qualità.
\paragraph{Rischi non previsti}
Il valore è identificato con un numero intero e indica il numero di rischi non preventivati.
Questi sono i valori da noi ritenuti opportuni:\\
\textit{Valore minimo}: 10\\
\textit{Valore ottimo}: 5\\
\subsection{Gestione Infrastrutture}
\subsubsection{GitHub\textsubscript{g}}
\label{GitHub} 
Servizio di hosting per progetti software che implementa uno strumento di controllo versione distribuito Git\textsubscript{g}.
Oltre alla copia in remoto del repository\textsubscript{g} di progetto ogni componente del gruppo ha una propria copia in locale.\\
Per ottenere una copia del repository\textsubscript{g} ogni componente ha scaricato lo strumento Git\textsubscript{g} ed eseguendo il 
comando 'git clone' da git\textsubscript{g} bash viene creata una cartella collegata alla repository\textsubscript{g} di progetto.\\
Non sono state imposte modalità specifiche sull'interazione con il repository\textsubscript{g} remoto in modo da non sconvolgere le abitudini di lavoro di 
ciascun componente.\\
I componenti del gruppo abituati ad interagire con GitHub\textsubscript{g} da interfaccia grafica possono continuare a farne uso.
\paragraph{Repository\textsubscript{g}}
Il repository\textsubscript{g} si può trovare all'indirizzo \textbf{\textit{https://github.com/7clickers/ShowRoom3D}} ed è pubblico. 
I collaboratori sono i componenti del gruppo SevenClickers che utilizzano il proprio account GitHub\textsubscript{g} personale per collaborare al progetto.
\paragraph{Branching}
I branch si dividono in: 
\begin{itemize}
	\item \textbf{Branches protetti}:
    \begin{itemize} 
        \item \textbf{main}: contiene le versioni di release\textsubscript{g} del software;
        \item \textbf{documentation}: contiene i template latex\textsubscript{g} e rispettivi pdf della documentazione. I documenti presenti 
		in documentation sono stati approvati dal responsabile o almeno verificati dai verificatori.
    \end{itemize}
	\item \textbf{Branches liberi}: viene creato un branch\textsubscript{g} libero per ogni documento che si sta modificando. Il nome del branch deve essere:
	\begin{itemize}	
		\item Uguale al nome del documento;
		\item Scritto tutto in minuscolo;
		\item Gli spazi tra le parole devono essere sostituiti con il trattino basso.
	\end{itemize}
	
	In ogni branch\textsubscript{g} libero possono lavorare uno o più componenti del gruppo, a condizione che ognuno lavori su parti diverse 
	del documento a cui fa riferimento il branch\textsubscript{g}. Ad ogni componente che lavora sul branch\textsubscript{g}  
	viene assegnata una issue da risolvere. Una volta che tutte le issue relative al branch\textsubscript{g} sono state risolte ed è avvenuta 
	la verifica, si fa il merge e il branch\textsubscript{g} viene eliminato.
\end{itemize}

\paragraph{Commits}
\label{Commits}
È preferibile che ogni commit abbia una singola responsabilità per cambiamento.\\
I commits non possono essere effettuati direttamente sui branch protetti ma per integrare delle aggiunte o modifiche sarà necessario aprire una Pull Request.
All'approvazione di una Pull Request tutti i commit relativi al merge verranno raggruppati in un unico commit che rispetti la struttura sintattica descritta in seguito.

I commit dovranno essere accompagnati da una descrizione solo se ritenuta indispensabile alla comprensione del commit stesso.

I messaggi di commit sui \textbf{\uppercase{branch protetti}} dovranno seguire la seguente struttura sintattica:\\

\textbf{$<$label$>$$<$\#n\_issue$>$$<$testo$>$}\\\\
dove:\\\\
\begin{itemize}
\item \textbf{label} può assumere i seguenti valori:
	\begin{itemize}
	\item \textbf{feat}: indica che è stata implementata una nuova funzionalità;
	\item \textbf{fix}: indica che è stato risolto un bug;
	\item \textbf{update}: indica che è stata apportata una modifica che non sia fix o feat;
	\item \textbf{test}: qualsiasi cosa legata ai test;
	\item \textbf{docs}: qualsiasi cosa legata alla documentazione.
	\end{itemize}
\item \textbf{n\_issue}: indica il numero della issue a cui fa riferimento il commit (se non fa riferimento a nessuna issue viene omesso);
\item \textbf{testo}: indica con quale branch è stato effettuato il merge e deve rispettare la forma: merge from $<$nome branch da integrare$>$ to $<$nome branch corrente$>$;
\item \textbf{descrizione}: se aggiunta ad un commit deve rispondere alle domande WHAT?, WHY?, HOW? ovvero
cosa è cambiato, perché sono stati fatti i cambiamenti,in che modo sono stati fatti i cambiamenti.\\
\end{itemize}
I messaggi di commit sui \textbf{\uppercase{branch liberi}} dovranno seguire la struttura sintattica dei branch protetti ad eccezione del testo.
Il testo dei commit sui branch liberi non è soggetto a restrizioni particolari a patto che indichi in maniera intuitiva i cambiamenti fatti
in modo che possano essere compresi anche dagli altri collaboratori.



\paragraph{Pull Requests}
\label{Pull_Requests}
Per effettuare un merge su un branch protetto si deve aprire da GitHub\textsubscript{g} una Pull Request.
La Pull Request permette di verificare il lavoro svolto prima di integrarlo con il branch desiderato.\\
Alla creazione di una Pull Request bisogna associare:
\begin{itemize}
\item I verificatori in carica hanno il compito di trovare eventuali errori o mancanze e fornire un feedback riguardante il contenuto direttamente su GitHub\textsubscript{g} richiedendo
una review con un review comment sulla parte specifica da revisionare o con un commento generico.\\
Non sarà possibile effettuare il merge finchè tutti i commenti di revisione non saranno stati risolti e la Pull Request approvata da due verificatori;
\item L’issue associata nell’opzione “Development” che verrà chiusa alla risoluzione della Pull Request;
\item La Projects Board di cui fa parte;
\item Gli assegnatari che hanno il compito di apportare le modifiche necessarie in fase di verifica;
\item Le labels associate.
\end{itemize}
Per i commit relativi alle Pull Requests seguire le regole descritte nella sezione \fullref{Commits} per i branch protetti.

\paragraph{Milestone\textsubscript{g}}
Una milestone indica un traguardo intermedio significativo per il progetto.
Ad essa possono venire assegnate delle issues per verificarne il raggiungimento.
Ogni milestone ha una scadenza che viene discussa e fissata da tutto il gruppo.
Oltre alle 2 milestone + 1 milestone falcoltativa fissate dal committente il gruppo ne creerá ulteriori per scandire piú nel dettaglio i passi che ci porteranno 
ad ottenere i risultati prefissati.
Una delle milestone create dal gruppo durante il completamento della tecnology baseline relativa alla prima milestone imposta dal committente (Requirements and Tecnology Baseline)
riguarda l'implementazione di un PoC (Proof of concept).

\paragraph{Projects Board\textsubscript{g}}
Vengono utilizzate due project board per tracciare le issues della repo.
Una project board principale utilizzata da tutti i membri del gruppo e una project board per le approvazioni utilizzata solo dal responsabile di progetto per approvare i file che richiedono l'approvazione prima di una consegna.

\begin{itemize}
	\item La projectboard\textsubscript{g} principale è suddivisa in queste sezioni:
	\begin{itemize}
		\item \textbf{Todo}: issue\textsubscript{g} che non sono ancora state iniziate o che non sono ancora state assegnate;
		\item \textbf{In Progress}: issue\textsubscript{g} che sono state assegnate e a cui almeno un membro a cui è stata assegnata ha iniziato a lavorarci;
		\item \textbf{Pull Request}: issue\textsubscript{g} che è in fase di integrazione e necessita della verifica dei verificatori. Corrisponde all'inizio di una pull request;
		\item \textbf{Done}: issue\textsubscript{g} che sono state chiuse e che sono state verificate (se necessitano di verifica\textsubscript{g});
		\item \textbf{Approved}: issue\textsubscript{g} con label "Da Approvare" che hanno ottenuto l'approvazione\textsubscript{g} del responsabile subito dopo la verifica;
		Per tutte le issues che richiedono un'approvazione solo al momento della consegna è stata creata la project board dedicata alle approvazioni.
	\end{itemize}
	Inoltre nella project board principale vengono registrate delle issue che non richiedono verifica, approvazione o neanche integrazione, con lo scopo di monitorare meglio il lavoro di ogni membro del team.\\
Queste issue verranno chiuse e archiviate manualmente una volta che avranno terminato la loro utilità, un esempio può essere la seguente issue:\\
\textbf{diario di bordo 21-11-22}; questa issue non necessità verifica, approvazione o integrazione perchè non è di interesse caricare il file nella repo, però è utile tracciare lo svolgimento della issue.

	\item La projectboard\textsubscript{g} riservata alle approvazioni è suddivisa in queste sezioni:
	\begin{itemize}
		\item \textbf{Todo}: issue\textsubscript{g} che non sono ancora state iniziate dal responsabile;
		\item \textbf{In Progress}: issue\textsubscript{g}  che sono state prese in carico per essere approvate dal responsabile;
		\item \textbf{Pull Request}: issue\textsubscript{g} che è in fase di integrazione e necessita della verifica dei verificatori. Corrisponde all'inizio di una pull request;
		in questo caso i verificatori dovranno solo controllare che l'intestazione e il registro delle modifiche siano stati compilati correttamente dal responsabile di progetto in quanto tutto il resto del contenuto è già stato verificato in precedenza;
		\item \textbf{Approved}: issue\textsubscript{g} che sono state approvate dal responsabile.
	\end{itemize}
\end{itemize}
Le issue che si trovano nella project board approvazioni sono state create come continuazione ad issues che hanno in precedenza attraversato il work flow della project board principale fino allo stato "Done".
In questo modo il Responsabile di progetto dovrà approvare solamente file che si trovano nei branches protetti.
Vengono ricreate in questa project board per  mantenerne la tracciabilità e permetterne una più facile reperibilità futura per l'approvazione.
Questo permette anche di pulire la project board principale da tutte quelle issues che rimarrebbero inattive per molto tempo.
\\\\
Esempio di work flow issues:
\begin{enumerate}
	\item Viene creata l'issue ed inserita all'interno della sezione "To Do" della project board principale;
	\item Quando viene presa in carico l'issue viene spostata nella sezione "In Progress" fino al suo completamento;
	\item L'incaricato alla risoluzione dell'issue apre una pull request a cui assegna l'issue in modo che  venga chiusa in automatico al momento dell'integrazione con il branch di destinazione;
	\item Vengono fatte le verifiche opportune dai verificatori e le conseguenti modifiche prima di accettare la pull request;
	\item Dopo l'accettazione la pull request viene postata insieme alle issues associate nella colonna "Done";
	\item A questo punto si possono presentare due casi: 
	\begin{enumerate}
		\item I file relativi alla pull request richiedono approvazione immediata.
		Viene archiviata la pull request (e le issue associate) e creata una issue a cui viene aggiunta la label "Da approvare" con il nome del file da approvare.
		Il responsabile creerà un branch per approvare le issues con approvazione immediata aprirà una pull request per integrare nel branch di destinazione l'approvazione;
		\item I file relativi alla pull request richiedono l'approvazione prima di una consegna.
		Il responsabile archivia la pull request e ne crea una issue con il nome dei file da approvare nella project board dedicata alle approvazioni.
		NB: se nella project board Approvazioni è già presente un'issue di approvazione per il file da approvare non viene creata una nuova issue di approvazione ma si tiene la precedente.
		L'issue segue il work flow della project board approvazioni.
	\end{enumerate}
\end{enumerate}
\paragraph{Issue Tracking System\textsubscript{g}}
Il gruppo utilizza l'issue tracking system\textsubscript{g} di Github\textsubscript{g} per tenere traccia delle issue\textsubscript{g}. 
Le issues\textsubscript{g} verranno determinate dal responsabile, ma la loro assegnazione verrà effettuata dai membri del gruppo, in base alla priorità delle issues\textsubscript{g}, i loro ruoli e alle loro disponibilità temporali.\\
Per marchiare le issues secondo criteri di interesse (come ambito o prioritá) vengono utilizzate le labels.
\\\\\textbf{Labels:}
\begin{itemize}
	\item \textbf{Da approvare}: indica una issue o pull request che necessita di approvazione;
	\item \textbf{Da verificare}: indica una issue o pull request che necessita di verifica;
	\item \textbf{Documentazione}: indica una issue o pull request riguardante la documentazione;
	\item \textbf{P=Bassa}: indica una issue o pull request a priorità bassa, con una scadenza di massimo 15 giorni;
	\item \textbf{P=Media}: indica una issue o pull request a priorità media, con una scadenza di massimo 7 giorni;
	\item \textbf{P=Alta}: indica una issue o pull request a priorità alta, con una scadenza di massimo 3 giorni;
	\item \textbf{Presentazione}: indica una issue riguardante la redazione di una presentazione in classe o al proponente.
\end{itemize}

Nel caso un membro del gruppo dovesse rendersi conto che l'issue\textsubscript{g} che sta svolgendo potrebbe essere suddiviso in ulteriori issue\textsubscript{g}, dovrà rivolgersi al responsabile, che è l'unico che può aggiungere, modificare o eliminare le issue\textsubscript{g}.\\
Per raggruppare più issue si marca ciascuna con la label di raggruppamento.
Il nome di una label di raggruppamento viene preceduto da una G maiuscola (che sta per gruppo) seguito dal nome dell'attività da svolgere. esempio: "G Piano di progetto"
Al completamento di tutte le issues di un gruppo si potrà scegliere se eliminare la label o tenerla per raggruppare issues future dello stesso gruppo.
indica un gruppo di labels relative al documento \textit{Piano di Progetto}.
\\\\
\textbf{Utilizzo delle checkbox:}
\\\\
Una issue può essere suddivisa in più attività tramite delle checkbox in base alla grandezza o
alla necessità di suddividere il lavoro. Al completamento di un'attività si spunta la checkbox
corrispondente in modo da avvisare il gruppo sullo stato di avanzamento di quella issue.

\subsubsection{Discord}
Strumento utilizzato per la comunicazione tra i componenti del gruppo. Sono creati diversi canali con diverse funzioni:
\begin{itemize}
	\item \textbf{Testuali}: si condividono informazioni in formato testuale;
	\item \textbf{Vocali}: si effettuano videochiamate brevi, che non necessitano di essere verbalizzate;
	\item \textbf{Risorse}: si condivide materiale da consultazione (ed esempio documenti, link, ecc...). 
\end{itemize}
I canali sono creati secondo le necessità del periodo, in accordo con tutti i componenti del gruppo.
Su Discord è anche possibile controllare in ogni momento quale ruolo è assegnato ad ogni membro del gruppo.


\pagebreak
\section{Standard di qualità ISO/IEC 9126}
Questa sezione descrive in dettaglio lo standard ISO/IEC 9126 utilizzato dal gruppo. Le norme di questo standard descrivono un modello di qualità del software, definiscono le caratteristiche che lo determinano e propongono metriche per la misurazione.
Le norme constano di quattro parti:
\begin{itemize}
\item Parte 1: Metriche per la qualità esterna
\item Parte 2: Metriche per la qualità interna
\item Parte 3: Modello della qualità del software
\item Parte 4: Metriche per la qualità in uso
\end{itemize}

\input{res/sections/isoiec_9126/qualità_esterne.tex}
\input{res/sections/isoiec_9126/qualità_interne.tex}
\input{res/sections/isoiec_9126/modello_qualità.tex}
\input{res/sections/isoiec_9126/qualità_in_uso.tex}

\pagebreak
\section{Standard di qualità ISO/IEC 12207:1995}
Questa sezione descrive in dettaglio lo standard ISO/IEC 12207:1995 scelto dal gruppo per garantire la qualità dei processi del ciclo di vita\textsubscript{g} di un software. I processi dello standard contengono attività e compiti che devono essere applicati durante l'acquisizione di un sistema software. Esso contiene tre tipologie di processi che sono: 
\begin{itemize}
\item Processi primari;
\item Processi di supporto;
\item Processi organizzativi.
\end{itemize}

\subsection{Processi primari}
Questi processi sono essenziali in quanto un progetto è detto tale se e solo se è attivo almeno un processo primario.
I processi di supporto definiti dallo standard sono i seguenti.
\subsubsection{Acquisizione}
Il processo di acquisizione contiene le attività ed i compiti dell'acquirente.
Le attività del processo sono le seguenti:
\begin{itemize}
\item Iniziazione
\item Preparazione della richiesta di proposta
\item Preparazione e aggiornamento del contratto
\item Monitoraggio dei fornitori
\item Accettazione e completamento
\end{itemize}

\subsubsection{Fornitura}
Il processo di fornitura contiene le attività e le mansioni del fornitore. Il processo può essere.
Le attività del processo sono le seguenti:
\begin{itemize}
\item Iniziazione
\item Preparazione alla risposta
\item Contratto
\item Pianificazione
\item Esecuzione e controllo
\item Revisione e valutazione
\item Rilascio e completamento
\end{itemize}

\subsubsection{Sviluppo}
Il processo di sviluppo contiene le attività e i compiti dello sviluppatore.
Il processo contiene le attività per l'analisi dei requisiti, la progettazione, la codifica, l'integrazione, il test, l'installazione e l'accettazione relative al sistema se previsto dal contratto.
Le attività del processo sono le seguenti:
\begin{itemize}
\item Implementazione dei processi
\item Analisi dei requisiti di sistema
\item Progettazione architettura di sistema
\item Analisi requisiti software
\item Progettazione architettura software
\item Progettazione dettagliata del software
\item Codifica e test del software
\item Integrazione software
\item Test di qualificazione del software
\item Integrazione del sistema 
\item Test di qualificazione del sistema 
\item Installazione software
\item Supporto per l'accettazione del software
\end{itemize}

\subsubsection{Gestione operativa}
Il processo di gestione operativa contiene le attività e i compiti dell'operatore. Il processo riguarda il funzionamento del prodotto software e il supporto operativo agli utenti. 
Le attività del processo sono le seguenti:
\begin{itemize}
\item Implementazione dei processi
\item Collaudo operativo
\item Operazione di sistema
\item Supporto all'utente
\end{itemize}

\subsubsection{Manutenzione}
Il processo di manutenzione contiene le attività e i compiti del manutentore. Questo processo viene attivato quando il prodotto software subisce modifiche al codice e alla documentazione associata a causa di un problema o necessità di miglioramento o adattamento.
Le attività del processo sono le seguenti:
\begin{itemize}
\item Implementazione di processi
\item Analisi dei problemi e delle modifiche
\item Implementanzione della modifica
\item Revisione/Accettazione della manuntenzione
\item Migrazione
\item Ritiro del software
\end{itemize}



\subsection{Processi di supporto}
I processi di supporto aiutano le attività di tutti gli altri processi dell'organizzazione a garantire il successo e la qualità del progetto. Questi processi possono essere attivati da un processo\textsubscript{g} primario o da un altro processo\textsubscript{g} di supporto.
Le attività e i compiti in un processo\textsubscript{g} di supporto sono di responsabilità dell'organizzazione che esegue tale processo\textsubscript{g}. Questa organizzazione garantisce che il processo\textsubscript{g} sia in atto e funzionante.
I processi di supporto definiti dallo standard sono i seguenti.
\subsubsection{Documentazione}
Il processo\textsubscript{g} di documentazione è un processo\textsubscript{g} per la registrazione delle informazioni prodotte da un processo\textsubscript{g} o attività del ciclo di vita\textsubscript{g}. Il processo\textsubscript{g} contiene l'insieme delle attività che pianificano, progettano, sviluppano, producono, modificano, distribuiscono e mantengono quei documenti necessari a tutti gli interessati.
Le attività del processo\textsubscript{g} sono le seguenti:
\begin{itemize}
\item Implementazione di processo\textsubscript{g};
\item Produzione;
\item Progettazione e sviluppo;
\item Manutenzione.
\end{itemize}
\subsubsection{Gestione della configurazione}
Il processo\textsubscript{g} di gestione della configurazione è un processo\textsubscript{g} di applicazione di procedure amministrative e tecniche durante tutto il ciclo di vita\textsubscript{g} del software per mantenere l'integrità di tutti i componenti della configurazione e di renderli accessibili a chi ne ha diritto.
Le attività del processo\textsubscript{g} sono le seguenti:
\begin{itemize}
\item Implementazione del processo\textsubscript{g};
\item Identificazione della configurazione;
\item Controllo della configurazione;
\item Valutazione dello stato di configurazione;
\item Gestione del rilascio e distribuzione.
\end{itemize}
\subsubsection{Accertamento della qualità}
Il processo\textsubscript{g} di accertamento della qualità è un processo\textsubscript{g} per fornire un'adeguata garanzia che i prodotti ed i processi software nel ciclo di vita\textsubscript{g} del progetto siano conformi ai requisiti specificati e aderiscano ai piani stabiliti. Per essere imparziale, la garanzia della qualità deve avere libertà organizzativa e autorità da parte delle persone direttamente responsabili dello sviluppo del prodotto\textsubscript{g} software o dell'esecuzione del processo\textsubscript{g} nel progetto.
Le attività del processo\textsubscript{g} sono le seguenti:
\begin{itemize}
\item Implementazione del processo\textsubscript{g};
\item Accertamento di prodotto\textsubscript{g};
\item Accertamento di processo\textsubscript{g};
\item Accertamento della qualità di sistema.
\end{itemize}
\subsubsection{Verifica\textsubscript{g}}
Il processo\textsubscript{g} di verifica\textsubscript{g} è un processo\textsubscript{g} per determinare se i prodotti software di un'attività soddisfano i requisiti o le condizioni loro imposti nelle attività precedenti.
Le attività del processo\textsubscript{g} sono le seguenti:
\begin{itemize}
\item Implementanzione del processo\textsubscript{g};
\item Verifica\textsubscript{g}.
\end{itemize}
\subsubsection{Validazione\textsubscript{g}}
Il processo\textsubscript{g} di validazione\textsubscript{g} è un processo\textsubscript{g} per determinare se i requisiti e il sistema finale soddisfano l'uso specifico previsto.
Le attività del processo\textsubscript{g} sono le seguenti:
\begin{itemize}
\item Implementazione del processo\textsubscript{g};
\item Validazione\textsubscript{g}.
\end{itemize}

\subsubsection{Revisione congiunta}
Il processo\textsubscript{g} di revisione congiunta è un processo\textsubscript{g} per valutare lo stato ed i prodotti di un'attività di un progetto in modo appropriato. Le revisioni congiunte avvengono sia a livello di gestione del progetto che a livello tecnico.
\begin{itemize}
\item Implementazione del processo\textsubscript{g};
\item Revisioni della gestione del progetto;
\item Revisioni tecniche.
\end{itemize}

\subsubsection{Audit}
Il processo\textsubscript{g} di Audit ha lo scopo di determinare in maniera indipendente la conformità di prodotti e processi selezionati ai requisiti, piani e accordi.
Le attività del processo\textsubscript{g} sono le seguenti:
\begin{itemize}
\item Implementazione del processo\textsubscript{g};
\item Audit.
\end{itemize}

\subsubsection{Risoluzione dei problemi}
Il Processo\textsubscript{g} di risoluzione dei problemi è un processo\textsubscript{g} per l'analisi e la risoluzione dei problemi, indipendentemente dalla loro natura o origine, che vengono scoperti durante l'esecuzione di processi di sviluppo, funzionamento, manutenzione o altri. L'obiettivo è fornire un mezzo tempestivo, responsabile e documentato per garantire che tutti i problemi scoperti siano analizzati e risolti e che le tendenze siano riconosciute.
\begin{itemize}
\item Implementazione del processo\textsubscript{g};
\item Risoluzione del problema.
\end{itemize}


\subsection{Processi organizzativi}
Le attività e i compiti in un processo\textsubscript{g} organizzativo sono di responsabilità dell'organizzazione che utilizza il processo\textsubscript{g}. L'organizzazione garantisce che il processo\textsubscript{g} sia attivo e funzionante.
I processi organizzativi definiti nello standard sono i seguenti.

\subsubsection{Gestione organizzativa}
Il processo\textsubscript{g} di gestione organizzativa ha lo scopo di organizzare e monitorare i processi per il raggiungimento dei loro obiettivi. Il processo\textsubscript{g} è stabilito da una organizzazione per assicurare la consistente applicazione di pratiche per l'uso dall'organizzazione e nei progetti.
Le attività del processo\textsubscript{g} sono le seguenti:
\begin{itemize}
\item Inizio e definizione dello scopo;
\item Pianificazione;
\item Esecuzione e controllo;
\item Revisione e valutazione;
\item Chiusura.
\end{itemize}
\subsubsection{Gestione delle infrastrutture}
Il processo\textsubscript{g} di gestione delle infrastrutture è un processo\textsubscript{g} per stabilire e mantenere l'infrastruttura necessaria per qualsiasi altro processo\textsubscript{g}. L'infrastruttura può includere hardware, software, strumenti, tecniche, standard e strutture per lo sviluppo, il funzionamento o la manutenzione.
Le attività del processo\textsubscript{g} sono le seguenti:
\begin{itemize}
\item Implementazione del processo\textsubscript{g};
\item Realizzazione dell'infrastruttura;
\item Manutenzione dell'infrastruttura.
\end{itemize}

\subsubsection{Miglioramento del processo\textsubscript{g}}
Il processo\textsubscript{g} di miglioramento è un processo\textsubscript{g} per stabilire, valutare, misurare, controllare e migliorare un processo\textsubscript{g} del ciclo di vita\textsubscript{g} del software.
Le attività del processo\textsubscript{g} sono le seguenti:
\begin{itemize}
\item Stabilimento del processo\textsubscript{g};
\item Valutazione del processo\textsubscript{g};
\item Miglioramento del processo\textsubscript{g}.
\end{itemize}

\subsubsection{Formazione del personale}
Il processo\textsubscript{g} di formazione del personale è un processo\textsubscript{g} per fornire e mantenere personale qualificato. L'acquisizione, la fornitura, lo sviluppo, il funzionamento o la manutenzione di prodotti software dipende in gran parte da personale esperto e qualificato.
Le attività del processo\textsubscript{g} sono le seguenti:
\begin{itemize}
\item Implementazione del processo\textsubscript{g};
\item Sviluppo materiale didattico;
\item Realizzazione piano formativo.
\end{itemize}
\pagebreak
\section{Metriche di qualità}

\subsection{Metriche per la qualità di processo}
In questa sezione sono descritte le metriche di qualità di processo che il gruppo intende adottare.

\subsubsection{MPC01: Planned Value (PV)}
Costo (in \euro) pianificato per realizzare le attività di progetto pianificate al momento del calcolo.
Per calcolare questa metrica si utilizzerà la formula:
\begin{equation*}
\text{PV}=\text{B}\textsubscript{tot} * \%\text{ lavoro pianificato}
\end{equation*}
\noindent Per eseguire il calcolo:
\begin{itemize}
\item la \% di lavoro pianificato viene calcolata dividendo le ore pianificate per le ore pianificate totali;
\item B\textsubscript{tot} indica il budget totale preventivato.
\end{itemize}
\textit{Valore minimo:} $\ge 0$ \euro\\
\textit{Valore ottimo:}	$\le \text{Budget at Completion\textsubscript{g}}$
\subsubsection{MPC02: Actual Cost (AC)}
Costo (in \euro) sostenuto al momento del calcolo.
Per calcolare questa metrica si tiene conto di:
\begin{equation*}
\text{AC}=\sum (\text{Totale ore lavorative per incremento}*\text{costo orario})
\end{equation*}
\noindent Per eseguire il calcolo:
\begin{itemize}
\item il costo orario è un costo fisso per ruolo ad ora;
\item Totale ore lavorative per incremento va a contare per incremento le ore effettive di lavoro per ruolo.
\end{itemize}
\textit{Valore minimo:} $\ge 0$ \euro\\
\textit{Valore ottimo:} $\le \text{EAC}$
\subsubsection{MPC03: Estimated at Completion (EAC)}
Un valore (in \euro) di revisione di costo stimato per la realizzazione del progetto alla data corrente.
\begin{equation*}
\text{EAC}=\text{AC}+\text{ETC}
\end{equation*}
\textit{Valore minimo:} 
$ \text{EAC} \le \text{preventivo} -8\% $ \textit{oppure} $ \text{EAC} \ge \text{preventivo} +5\% $\\
\textit{Valore ottimo:} Costo preventivato
\subsubsection{MPC04: Earned Value (EV)}
Indica il valore (in \euro) delle attività realizzate nel progetto fino al momento della misurazione.
\begin{equation*}
\text{EV}=\text{B\textsubscript{tot} * \%\text{ lavoro effettivo}}
\end{equation*}
\noindent Per eseguire il calcolo:
\begin{itemize}
\item la \% lavoro effettivo è calcolata dividendo le ore effettive per le ore totali preventivate.
\item B\textsubscript{tot} indica il budget totale preventivato.
\end{itemize}
\textit{Valore minimo:} $\ge 0$ \euro\\
\textit{Valore ottimo:} $\le \text{EAC}$
\subsubsection{MPC05: Estimated to Complete (ETC)}
Metrica che produce un valore (in \euro) per indicare il numero di attività necessarie al completamento del progetto.\\
\begin{equation*}
\text{ETC}=B\textsubscript{tot}-\text{EV}
\end{equation*}
\noindent Per eseguire il calcolo:
\begin{itemize}
\item B\textsubscript{tot} indica il budget totale preventivato.
\end{itemize}
\textit{Valore minimo:} $\ge 0$ \euro\\
\textit{Valore ottimo:} $\le \text{EAC}$
\subsubsection{MPC06: Cost Variance (CV)}
Valore (in \euro)che indica se il valore del costo attuale del progetto è maggiore, uguale o minore rispetto al costo effettivo. Se CV $\ge 0$ significa che si sta producendo con maggior efficienza risparmiando sul costo rispetto a quanto pianificato, viceversa se negativo.
\begin{equation*}
\text{CV}=\text{EV}-\text{AC}
\end{equation*}
\textit{Valore minimo:} CV = 0 \euro\\
\textit{Valore ottimo:} CV $\ge 0$ \euro
\subsubsection{MPC07: Schedule Variance (SV)}
Variazione percentuale utilizzata per analizzare se il progetto è in linea, in anticipo o in ritardo rispetto alla schedulazione delle attività di progetto pianificate nella baseline.
\begin{equation*}
\text{SV}=\frac{\text{EV}-\text{PV}}{\text{EV}}*100
\end{equation*}
\textit{Valore minimo:} $\ge -15\%$\\
\textit{Valore ottimo:} 0\%
\subsubsection{MPC08: Budget Variance (BV)}
Valore (in \euro) che indica se alla data corrente si è speso di più o di meno del budget iniziale. Un valore di BV $\ge 0$ significa che il progetto sta spendendo il proprio budget con minor velocità di quanto pianificato, viceversa se negativo.
\begin{equation*}
\text{BV}=\text{PV}-\text{AC}
\end{equation*}
\textit{Valore minimo:} BV = 0 \euro\\
\textit{Valore ottimo:} BV $\ge 0 $ \euro

\subsubsection{MPC09: Requirements Stability Index (RSI)}
Indice che traccia la variazione dei requisiti nell'arco del progetto.\\
Si calcola tramite la seguente formula:
\begin{equation*}
\text{RSI}=1-\frac{\text{numero requisiti cambiati}+\text{numero requisiti eliminati}+\text{numero requisiti aggiunti}}{\text{numero totale dei requisiti iniziali}}*100
\end{equation*}
\textit{Valore minimo:} $70\%$\\
\textit{Valore ottimo:} $100\%$

\subsubsection{MPC10: Indice di Gulpease}
Si tratta dell'indice di leggibilità di un testo tarato sulla lingua italiana.
I risultati sono compresi tra 0 e 100, dove il valore "100" indica la leggibilità più alta e "0" la leggibilità più bassa. Ai seguenti valori si associano i seguenti significati:
\begin{itemize}
\item inferiore a 80 sono difficili da leggere per chi ha la licenza elementare
\item inferiore a 60 sono difficili da leggere per chi ha la licenza media
\item inferiore a 40 sono difficili da leggere per chi ha un diploma superiore
\end{itemize}
Viene adottata la seguente formula per calcolarlo:
\begin{equation*}
89+\frac{300*(\text{numero delle frasi})-10*(\text{numero delle lettere})}{\text{numero delle parole}}
\end{equation*}\\
Questi sono i valori da noi ritenuti opportuni:\\
\textit{Valore minimo}: $ \ge 50 $\\ 
\textit{Valore ottimo}: $ \ge 80 $\\

\subsubsection{MPC11: Metriche soddisfatte}
Viene calcolato questo indice tramite una percentuale per tenere conto delle metriche soddisfatte. Una metrica si dice "soddisfatta" se raggiunge almeno il valore minimo imposto sul \textit{Piano di qualifica}.\\
Questi sono i valori da noi ritenuti opportuni:\\\\
\textit{Valore minimo:} $ \ge 80\% $\\
\textit{Valore ottimo:} $ 100\% $\\

\subsubsection{MPC12: Code coverage}
Viene definita come la percentuale di linee di codice del progetto che sono state eseguite dai test dopo un'esecuzione.\\
Questi sono i valori da noi ritenuti opportuni:\\
\textit{Valore minimo}: $ \ge 70\% $\\
\textit{Valore ottimo}: $ \ge 90\% - 100\% $\\

\subsubsection{MPC13: Rischi non previsti}
Il valore è identificato con un numero intero e indica il numero di rischi non previsti durante il corso del progetto.\\
Questi sono i valori da noi ritenuti opportuni:\\
\textit{Valore minimo}: $ \ge 0$\\
\textit{Valore ottimo}: 0\\
\subsection{Metriche per la qualità di prodotto\textsubscript{g}}
In questa sezione sono descritte le metriche di qualità di prodotto\textsubscript{g} che il gruppo intende adottare.

\subsubsection{MPD01: Percentuale requisiti soddisfatti}
Valore percentuale che serve per indicare i requisiti soddisfatti.\\
\textit{Valore minimo:} $100\% \text{ dei requisiti obbligatori}$\\
\textit{Valore ottimo:} $100\% \text{ di tutti i requisiti}$

\subsubsection{MPD02: Densità di fallimenti durante l'esecuzione}
Si intende la percentuale di failure\textsubscript{g} o di esecuzioni non andate a buon fine di determinate azioni.
 Le eventuali esecuzioni fallite o failure\textsubscript{g} sono state segnate dai Programmatori con cui hanno poi calcolato il valore della metrica.
\begin{equation*}
\text{Densità di fallimenti}=\frac{\text{numero test eseguiti sul programma falliti}}{\text{numero totale di test eseguiti}}*100
\end{equation*}
\textit{Valore minimo:} $20\%$\\
\textit{Valore ottimo:} $10\%$

\subsubsection{MPD03: Tempo medio di risposta}
Tempo medio impiegato dal software per rispondere a una richiesta utente o svolgere un’attività di sistema.
\textit{Valore minimo}: 4 secondi\\
\textit{Valore ottimo}: 2 secondi\\

\subsubsection{MPD04: Tempo di caricamento}
Tempo medio di attesa per il caricamento del sito.\\
\textit{Valore minimo}: 15 secondi\\
\textit{Valore ottimo}: 10 secondi\\

\subsubsection{MPD05: Facilità di apprendimento}
Misura l'intuibilità e la facilità di utilizzo del programma.\\
\textit{Valore minimo}: 5 minuti\\
\textit{Valore ottimo}: 2 minuti\\

\subsubsection{MPD06: Complessità ciclomatica}
Metrica utilizzata per misurare la complessità di un programma. Calcolata sul grafo dei cammini linearmente indipendenti percorsi dal software ed i nodi presenti, cioè i punti decisionali del programma.\\
\begin{equation*}
v(G) = e - n + 2p
\end{equation*}
dove:
\begin{itemize}
	\item \textit{v(G)} = complessità ciclomatica del grafo G
	\item \textit{e} = il numero di archi nel grafo
	\item \textit{n} = il numero di nodi nel grafo
	\item \textit{p} = il numero di componenti connesse
\end{itemize}
\textit{Valore minimo}: $ \le 10 $\\
\textit{Valore ottimo}: $ \le 4 $\\

\subsubsection{MPD07: Densità dei commenti}
Misura la percentuale delle righe di commento sul totale delle righe di codice presenti in un modulo.\\
\begin{equation*}
\text{Densità dei commenti}=\frac{\text{numero righe di commento}}{\text{numero righe di codice}}*100
\end{equation*}
\textit{Valore minimo}: 20\%\\
\textit{Valore ottimo}: 10\%\\

\subsubsection{MPD08: Browser supportati}
Valore percentuale dei browser supportati dal prodotto\textsubscript{g} software.\\
Da associare a questa metrica è doveroso inserire la versione del browser da cui è stato testato il prodotto\textsubscript{g}.
\begin{equation*}
\text{Browser supportati}=\frac{\text{browser in cui il prodotto\textsubscript{g} funziona}}{\text{browser testati}}*100
\end{equation*}	
\textit{Valore minimo}: 80\%\\
\textit{Valore ottimo}: 100\%\\

\pagebreak

\end{document}
